%%
%% This is file `sigconf.tex',
%% generated with the docstrip utility.
%%
%% The original source files were:
%%
%% samples.dtx  (with options: `all,proceedings,sigconf')
%% 
%% IMPORTANT NOTICE:
%% 
%% For the copyright see the source file.
%% 
%% Any modified versions of this file must be renamed
%% with new filenames distinct from sigconf.tex.
%% 
%% For distribution of the original source see the terms
%% for copying and modification in the file samples.dtx.
%% 
%% This generated file may be distributed as long as the
%% original source files, as listed above, are part of the
%% same distribution. (The sources need not necessarily be
%% in the same archive or directory.)
%%
%%
%% Commands for TeXCount
%TC:macro \cite [option:text,text]
%TC:macro \citep [option:text,text]
%TC:macro \citet [option:text,text]
%TC:envir table 0 1
%TC:envir table* 0 1
%TC:envir tabular [ignore] word
%TC:envir displaymath 0 word
%TC:envir math 0 word
%TC:envir comment 0 0
%%
%% The first command in your LaTeX source must be the \documentclass
%% command.
%%
%% For submission and review of your manuscript please change the
%% command to \documentclass[manuscript, screen, review]{acmart}.
%%
%% When submitting camera ready or to TAPS, please change the command
%% to \documentclass[sigconf]{acmart} or whichever template is required
%% for your publication.
%%
%%
\documentclass[sigconf]{acmart}
%%
%% \BibTeX command to typeset BibTeX logo in the docs
\AtBeginDocument{%
  \providecommand\BibTeX{{%
    Bib\TeX}}}

%% Rights management information.  This information is sent to you
%% when you complete the rights form.  These commands have SAMPLE
%% values in them; it is your responsibility as an author to replace
%% the commands and values with those provided to you when you
%% complete the rights form.
\setcopyright{acmlicensed}
\copyrightyear{2026}
\acmYear{2026}
\acmDOI{XXXXXXX.XXXXXXX}
%% These commands are for a PROCEEDINGS abstract or paper.
% \acmConference[ICCAD '26]{International Conference on Computer-Aided Design}{November 08--12, 2026}{San Jose, CA, USA}
%%
%%  Uncomment \acmBooktitle if the title of the proceedings is different
%%  from ``Proceedings of ...''!
%%
%%\acmBooktitle{Woodstock '18: ACM Symposium on Neural Gaze Detection,
%%  June 03--05, 2026, Woodstock, NY}
\acmISBN{978-1-4503-XXXX-X/2026/06}




%%
%% Submission ID.
%% Use this when submitting an article to a sponsored event. You'll
%% receive a unique submission ID from the organizers
%% of the event, and this ID should be used as the parameter to this command.
%%\acmSubmissionID{123-A56-BU3}

%%
%% For managing citations, it is recommended to use bibliography
%% files in BibTeX format.
%%
%% You can then either use BibTeX with the ACM-Reference-Format style,
%% or BibLaTeX with the acmnumeric or acmauthoryear sytles, that include
%% support for advanced citation of software artefact from the
%% biblatex-software package, also separately available on CTAN.
%%
%% Look at the sample-*-biblatex.tex files for templates showcasing
%% the biblatex styles.
%%
%%
%% The majority of ACM publications use numbered citations and
%% references.  The command \citestyle{authoryear} switches to the
%% "author year" style.
%%
%% If you are preparing content for an event
%% sponsored by ACM SIGGRAPH, you must use the "author year" style of
%% citations and references.
%% Uncommenting
%% the next command will enable that style.
%%\citestyle{acmauthoryear}


%%
\usepackage{braket}
\usepackage{xcolor}
\usepackage[table]{xcolor}
\usepackage{amsthm}
\usepackage{graphicx}
\usepackage{subcaption}
\usepackage{forest}
\usepackage{tikz}
\usetikzlibrary{positioning,arrows.meta,calc}
\newtheorem{proposition}{Proposition}

% Uncomment to show comments
% \def\SHOWCOMMENTS{}

\ifdefined\SHOWCOMMENTS

    \newcommand{\anik}[1]{%
        \textcolor{blue}{\textbf{Anik:} #1}%
    }

    \newcommand{\suman}[1]{%
        \textcolor{red}{\textbf{Suman:} #1}%
    }

    \newcommand{\siyi}[1]{%
        \textcolor{cyan}{\textbf{Siyi:} #1}%
    }

\else

    \newcommand{\anik}[1]{}
    \newcommand{\suman}[1]{}
    \newcommand{\siyi}[1]{}

\fi



\begin{document}

%%
%% The "title" command has an optional parameter,
%% allowing the author to define a "short title" to be used in page headers.
% \subtitle{Invited Paper}

% \title{\normalsize Invited Paper\\[0.3em]
%        \Huge Structure-Aware Placement and Routing of Multi-Controlled Toffoli on Bivariate Bicycle Code Architectures}
\title{%
{\normalsize Invited Paper}\\
Structure-Aware Placement and Routing of Multi-Controlled Toffoli on Bivariate Bicycle Code Architectures}
%%
%% The "author" command and its associated commands are used to define
%% the authors and their affiliations.
%% Of note is the shared affiliation of the first two authors, and the
%% "authornote" and "authornotemark" commands
%% used to denote shared contribution to the research.
% \author{Ben Trovato}
% \authornote{Both authors contributed equally to this research.}
% \email{trovato@corporation.com}
% \orcid{1234-5678-9012}
% \author{G.K.M. Tobin}
% \correspondingauthor
% \authornotemark[1]
% \email{webmaster@marysville-ohio.com}
% \affiliation{%
%   \institution{Institute for Clarity in Documentation}
%   \city{Dublin}
%   \state{Ohio}
%   \country{USA}
% }

%\author{Anik Basu Bhaumik}
%\authornotemark[1]
%\email{anikbasu001@e.ntu.edu.sg}
%\author{Suman Dutta}
%\authornotemark[1]
% \email{sumand.iiserb@gmail.com}
% \author{Siyi Wang}
% \email{siyi002@e.ntu.edu.sg}
% \author{Anupam Chattopadhyay}
% \email{anupam@ntu.edu.sg}
% \affiliation{%
%   \institution{Nanyang Technological University}
%   \country{Singapore}}


\author{Anik Basu Bhaumik}
\orcid{0009-0000-0778-2234}
\affiliation{%
  \institution{Nanyang Technological University}
  %\city{Singapore}
  \country{Singapore}
}
\email{anikbasu001@e.ntu.edu.sg}

\author{Suman Dutta}
\orcid{0009-0001-4891-9919}
\affiliation{%
\institution{Nanyang Technological University}
 %\city{Singapore}
 \country{Singapore}
 }
 \email{sumand.iiserb@gmail.com}

\author{Siyi Wang}
\orcid{0009-0006-9128-1857}
\affiliation{%
  \institution{Nanyang Technological University}
  %\city{Singapore}
  \country{Singapore}
  }
  \email{siyi002@e.ntu.edu.sg}

\author{Anupam Chattopadhyay}
\orcid{0000-0002-8818-6983}
\affiliation{%
  \institution{Nanyang Technological University}
  %\city{Singapore}
  \country{Singapore}
  }
\email{anupam@ntu.edu.sg}

\renewcommand{\shortauthors}{Basu Bhaumik et al.}


%%%%%%
\begin{abstract}
% The multi-controlled Toffoli (MCT) gate is a fundamental primitive in quantum circuit design, with broad applications in quantum arithmetic, cryptanalysis and algorithmic implementation. MCT being a high-level logical gate, decomposing MCT gates into lower-level netlists becomes of utmost importance and in dire need of optimization. While recent error-correcting codes such as the bivariate bicycle (BB) codes offer the potential for substantially reduced fault tolerance overhead, realizing non-Clifford circuits on modular BB-code architectures introduces a new compilation challenge where the cost of a logical operation depends upon the inter-module connections, factory density and placement at the code-level. Despite their importance, the mapping of multi-controlled Toffoli gates onto BB-code architectures remains relatively unexplored. In this paper, we map optimal-Toffoli-depth MCT decompositions (Dutta et al., PRA, 2025) onto BB-code-based fault-tolerant architectures. Each Toffoli gate is implemented through direct $\ket{\mathrm{CCZ}}$ state injection from an external magic state factory. We introduce a placement strategy that uses the binary-tree structure of the MCT decomposition algorithm to co-locate interacting subtrees and eventually construct larger MCTs, reducing inter-module instruction counts by up to $\mathbf{16.02}\%$ relative to a naive sequential first-fit placement strategy. In a comparison of linear and grid-like architectures under varying factory densities and placements, we observe up to a $\mathbf{23.7}\%$ reduction in inter-module instructions for grid-based layouts relative to linear architectures (Yoder et al., arXiv, 2025). In our analysis, we also present our circuits with their aggregate execution error and logical failure probability, evaluated from the Qiskit community bicyle-ISA error estimator \texttt{bicyle\_numerics} 

The multi-controlled Toffoli (MCT) gate is a fundamental primitive in quantum circuit design, with applications in quantum arithmetic, cryptanalysis, and algorithmic implementations. Being a high-level logical operation, the efficient decomposition of MCT gates into lower-level netlists has remained a major optimization challenge for decades. While emerging quantum error-correcting codes such as bivariate bicycle (BB) codes drastically reduce fault-tolerance overhead, realizing non-Clifford circuits on modular BB-code architectures introduces complex compilation bottlenecks governed by inter-module routing, factory density, and layout. Consequently, the mapping of MCT gates onto BB-code architectures remains relatively unexplored.

In this paper, we overcome these challenges by mapping optimal-Toffoli-depth MCT decompositions (Dutta et al., PRA, 2025) onto BB-code-based fault-tolerant architectures via direct $\ket{\mathrm{CCZ}}$ state injection from an external magic state factory. We introduce a targeted placement strategy that exploits the binary-tree structure of MCT decompositions to co-locate interacting subtrees. This approach reduces inter-module instruction counts by up to $\mathbf{16.02}\%$ compared to a naive sequential first-fit placement. We also evaluate the impact of factory placement across different topologies, demonstrating that grid-based layouts yield up to a $\mathbf{23.7}\%$ reduction in inter-module instructions relative to linear architectures (Yoder et al., arXiv, 2025). Finally, we validate the practical viability of our compiled circuits by analyzing aggregate execution errors and logical failure probabilities using the bicycle-ISA error estimator \texttt{bicycle\_numerics} provided by the Qiskit community\footnote{\url{https://github.com/qiskit-community/bicycle-architecture-compiler}}.
% \siyi{seems a bit too long. But this is not a big problem.}
% efficient implementations of MCT gates on such architectures have become increasingly important. In this paper, we map optimal-Toffoli-depth MCT decompositions (Dutta et al., PRA, 2025) onto BB-code-based fault-tolerant architectures. We implement each Toffoli gate via direct $\ket{\mathrm{CCZ}}$-state injection from an external magic-state factory, resulting in non-local Pauli-product measurements whose communication cost depends on logical-qubit placement and the density and placement of magic-state factories. We introduce a structure-aware placement strategy that exploits the binary-tree structure of the MCT decomposition to co-locate smaller MCTs and efficiently construct larger ones, showing upto $16.02\%$ improvement in inter-module instruction counts over a first-fit sequential placement strategy. We further evaluate the resulting mappings on linear and grid-like BB-module topologies under varying factory densities and placements, showing up to 23.7\% reduction in inter-module instruction counts over linear architectures. In our analysis, we present our Toffoli benchmark syntheses with the resulting aggregate error and logical failure probability, evaluated from the estimator \texttt{bicyle\_numerics} as provided by the Qiskit-community bicycle architecture compiler\footnote{\url{https://github.com/qiskit-community/bicycle-architecture-compiler}}.

\end{abstract}

%%
%% The code below is generated by the tool at http://dl.acm.org/ccs.cfm.
%% Please copy and paste the code instead of the example below.
%%
\begin{CCSXML}
<ccs2012>
 <concept>
  <concept_id>00000000.0000000.0000000</concept_id>
  <concept_desc>Do Not Use This Code, Generate the Correct Terms for Your Paper</concept_desc>
  <concept_significance>500</concept_significance>
 </concept>
 <concept>
  <concept_id>00000000.00000000.00000000</concept_id>
  <concept_desc>Do Not Use This Code, Generate the Correct Terms for Your Paper</concept_desc>
  <concept_significance>300</concept_significance>
 </concept>
 <concept>
  <concept_id>00000000.00000000.00000000</concept_id>
  <concept_desc>Do Not Use This Code, Generate the Correct Terms for Your Paper</concept_desc>
  <concept_significance>100</concept_significance>
 </concept>
 <concept>
  <concept_id>00000000.00000000.00000000</concept_id>
  <concept_desc>Do Not Use This Code, Generate the Correct Terms for Your Paper</concept_desc>
  <concept_significance>100</concept_significance>
 </concept>
</ccs2012>
\end{CCSXML}

\ccsdesc[500]{Bivariate Bicycle (BB) Code}
\ccsdesc{Fault-Tolerant Quantum Computing (FTQC)}
\ccsdesc[300]{Multi-Controlled Toffoli (MCT) Gate}
\ccsdesc{Quantum Circuit Design}
\ccsdesc[100]{Quantum Error-Correcting Codes}

%% A "teaser" image appears between the author and affiliation
%% information and the body of the document, and typically spans the
%% page.


\received{xx xxx 2026}
\received[revised]{xx xxx 2026}
\received[accepted]{xx xxx 2026}

%%
%% This command processes the author and affiliation and title
%% information and builds the first part of the formatted document.
\maketitle
% \vspace{-2em}
% Introduction
\section{Introduction}
\label{sec:intro}
Quantum gates are the fundamental building blocks of quantum circuits~\cite{gates1995}. Unlike classical gates, quantum gates are inherently reversible and are mathematically represented by unitary matrices. Among them, the multi-controlled Toffoli (MCT) gate is the most significant one with widespread applications in arithmetic circuit design~\cite{arithmetic}, reversible computation, oracle construction~\cite{dutta2026qic}, and quantum error-correction subroutines. Since MCT gates are logical quantum operations, they must be decomposed into a universal gate set~\cite{universal-gate2005}, such as the Clifford+T gate set. Consequently, efficient decomposition of MCT gates into the Clifford+T gate set has received considerable attention over the past two decades. Existing approaches have optimized various logical-level metrics, including T-count~\cite{gidney2021cccz}, T-depth~\cite{dutta2025pra}, ancilla count~\cite{nie2024arxiv}, and overall gate complexity, yielding several asymptotically optimal constructions under unrestricted~\cite{khattar2025quantum,dutta2025pra} and restricted~\cite{bhaumik2026optimal} qubit connectivity, assuming ideal, error-free execution. However, the efficient implementation of MCT gates in an error-corrected architecture and their resource constraints remain relatively underexplored in the literature.

Quantum error correction~\cite{shor-code,qecc-25years} is widely regarded as an enabling technology for large-scale quantum computation. Although noisy intermediate-scale quantum (NISQ)~\cite{preskill2018quantum} devices have demonstrated remarkable experimental progress, their limited coherence times and high physical error rates prevent the reliable execution of deep quantum circuits. Fault-tolerant quantum computing~\cite{ftqc} overcomes these limitations by encoding logical qubits into quantum error-correcting codes (QECCs), thereby enabling arbitrarily long computations provided that the physical error rate remains below the fault-tolerance threshold. Consequently, the practical cost of executing a quantum algorithm is determined not only by its logical circuit complexity, but also by the substantial overhead introduced by fault-tolerant error correction. Accordingly, recent resource-estimation studies increasingly account for these error-correction costs when evaluating the feasibility of large-scale quantum algorithms~\cite{gidney2025factor,cain2026shor}.

In an error-corrected architecture, every logical operation must be performed on encoded qubits, which internally involve syndrome extraction, logical gate execution, and decoder-assisted error recovery. As a result, the physical resources required to implement an error-corrected MCT circuit depend not only on its logical decomposition but also on the underlying quantum error-correcting code, the realization of logical Clifford+T gates, decoder complexity, and hardware communication constraints. Therefore, logical optimization metrics such as T-count and T-depth alone are insufficient to capture the practical cost of large-scale MCT implementations. In this work, we aim to bridge this gap by evaluating these implementations using BB-code-specific resource metrics.

The surface code~\cite{surface-code} has long been a leading candidate for fault-tolerant quantum computation owing to its high error threshold and compatibility with local interactions. Nevertheless, its substantial physical-qubit overhead has motivated the development of quantum low-density parity-check (qLDPC) codes~\cite{qLDPC}, which offer asymptotically higher encoding rates while preserving sparse stabilizer measurements. Among these, BB codes are the most promising for fault-tolerant architectures due to their structured connectivity and favorable code-implementation trade-offs. These developments motivate a re-evaluation of quantum resource estimation for MCT gates in the qLDPC-based architectures.


% Quantum error correction~\cite{shor-code,qecc-25years} is widely regarded as an enabling technology for large-scale quantum computation. Although noisy intermediate-scale quantum (NISQ)~\cite{preskill2018quantum} devices have demonstrated remarkable experimental progress, their limited coherence times and high physical error rates prevent the reliable execution of deep quantum circuits. Fault-tolerant quantum computing~\cite{ftqc} overcomes these limitations by encoding logical qubits into quantum error-correcting codes (QECCs), thereby enabling arbitrarily \anik{need a small mention of FTQC term in this para}long computations provided that the physical error rate remains below the fault-tolerance threshold. Consequently, the practical cost of executing a quantum algorithm is determined not only by its logical circuit complexity, but also by the substantial overhead introduced by fault-tolerant error correction. Accordingly, recent resource-estimation studies increasingly account for these error-correction costs when evaluating the feasibility of large-scale quantum algorithms~\cite{gidney2025factor,cain2026shor}.
% \siyi{Ok so the logic is not complete here. 1st para: QEC is importamnt 2nd para: MCT is significant q operation. lack of a highlight sentence for define this problem: QEC on MCT is very important. And then you can say the limitations and under exploration blabla as what u wrote in para 3.}

% % Quantum gates are the fundamental building blocks of quantum circuits~\cite{gates1995}. Unlike classical gates, quantum gates are inherently reversible and are mathematically represented by unitary matrices. Among them, the multi-controlled Toffoli (MCT) gate is the most significant one with widespread applications in arithmetic circuit design~\cite{arithmetic}, reversible computation, oracle construction~\cite{dutta2026qic}, and quantum error-correction subroutines. Since MCT gates are complex quantum operations, they must be decomposed into a universal gate set~\cite{universal-gate2005}, such as the Clifford+T gate set. Consequently, efficient decomposition of MCT gates into the Clifford+T gate set has received considerable attention over the past two decades. Existing approaches have optimized various logical-level metrics, including T-count~\cite{gidney2021cccz}, T-depth~\cite{dutta2025pra}, ancilla count~\cite{nie2024arxiv}, and overall gate complexity, yielding several asymptotically optimal constructions under unrestricted~\cite{khattar2025quantum,dutta2025pra} and restricted~\cite{bhaumik2026optimal} qubit connectivity, assuming ideal, error-free execution.
% The multi-controlled Toffoli (MCT) gate is a significant logical gate in quantum computation. This gate is a reversible representation of the classical multi-input AND gate. Since MCT gates are logical quantum operations, their physical realisations must be realised into a universal netlist gate set~\cite{universal-gate2005}, such as the Clifford+T gate set. Consequently, efficient decomposition of MCT gates into the Clifford+T gate set has received considerable attention over the past two decades. Existing approaches have optimized various logical-level metrics, including T-count~\cite{gidney2021cccz}, T-depth~\cite{dutta2025pra}, ancilla count~\cite{nie2024arxiv}, and overall gate complexity, yielding several asymptotically optimal constructions under unrestricted~\cite{khattar2025quantum,dutta2025pra} and restricted~\cite{bhaumik2026optimal} qubit connectivity, assuming ideal, error-free execution.
 

% However, considering an error-corrected architecture Toffoli related estimations and compilations still are relatively under-explored. In this paradigm every logical operation must be performed on encoded qubits and these qubits internally involves syndrome extraction, logical gate execution, and decoder-assisted error recovery. As a result, the physical resources required to implement an MCT circuit depend not only on its logical decomposition but also on the underlying quantum error-correcting code, the realization of logical Clifford+T gates, decoder complexity, and hardware communication constraints. Therefore, logical optimization metrics such as T-count and T-depth alone are insufficient to capture the practical cost of large-scale MCT implementations. In this work, we aim to bridge this gap by evaluating these implementations using BB-code-specific resource metrics.

% The surface code~\cite{surface-code} has long been a leading candidate for fault-tolerant quantum computation owing to its high error threshold and compatibility with local interactions. Nevertheless, its substantial physical-qubit overhead has motivated the development of quantum low-density parity-check (qLDPC) codes~\cite{qLDPC}, which offer asymptotically higher encoding rates while preserving sparse stabilizer measurements. Among these, BB codes are promising for fault-tolerant architectures due to structured connectivity and favorable code–implementation trade-offs. These developments motivate a re-evaluation of quantum resource estimation for MCT gates in qLDPC-based architectures
% \siyi{Ok but intro should be a overall perspective: for example now u are saying surface code->BB; but actually should be various QECC->Surface code, but -->BB}
% Recent advances in qLDPC code constructions, such as toric codes~\cite{kitaev2003}, bivariate bicycle (BB) codes~\cite{bravyi2024highthreshold}, quantum Tanner codes~\cite{leverrier2022focs}, balanced-product codes~\cite{bpqc2021}, and lifted-product codes~\cite{lpqc}, have established constant-rate code families with linear distance, significantly strengthening the prospects for scalable fault-tolerant quantum computing. .

% Among these constructions, BB codes are particularly promising for fault-tolerant implementations due to their structured stabilizer connectivity and favorable trade-offs between code parameters and implementation complexity. In tis work we aim to study the mapping of Toffoli gates onto this BB-module based architecture

%Some of the most popular varieties of qLDPC codes are listed as follows.\\
% Quantum LDPC (qLDPC) Codes:
% \begin{itemize}
%     \item Topological \& Geometrically Local
%     \begin{enumerate}
%         \item Surface / Toric Codes
%         \item Hyperbolic Surface Codes
%     \end{enumerate}
%     \item Homological \& Algebraic Products
%     \begin{enumerate}
%         \item Bivariate Bicycle (BB) Codes
%         \item Balanced Product Codes
%         %\item %Fiber Bundle Codes
%         \item Hypergraph Product (HGP) Codes
%     \end{enumerate}
%     \item Asymptotically Good Codes:
%     \begin{enumerate}
%         \item Lifted Product Codes
%         \item Generalized Tanner Codes
%     \end{enumerate}
%     \item Dynamic \& Subsystem Codes
%     \begin{enumerate}
%         \item Floquet / Honeycomb Codes
%     \end{enumerate}
% \end{itemize}
% In this paper, we investigate the fault-tolerant implementation of both basic and multi-controlled Toffoli gates under qLDPC-based quantum error correction. Rather than proposing new logical decompositions, we develop a qLDPC-aware resource-estimation framework to quantify the physical cost of state-of-the-art Toffoli and MCT decompositions after incorporating quantum error correction. We analyze the impact of encoding overhead, logical gate implementation, syndrome extraction, decoder complexity, and magic-state consumption on the overall resource requirements. The proposed framework enables a systematic comparison of different decomposition strategies using practical fault-tolerant metrics and provides insights into the trade-offs between logical optimization and physical implementation cost.

\subsection{Organization and contribution}
\label{sub:org}
% \anik{highlight the problem}
In this work, we address the lack of fault-tolerant resource estimates for Toffoli mapping by evaluating multi-controlled Toffoli (MCT) gates on bivariate bicycle (BB) code architectures. We decompose MCT gates into Toffoli gates and develop mapping techniques that exploit the decomposition's binary-tree structure. Our main contributions are as follows:
% \anik{elaborate a bit}
\begin{itemize}
    \item We establish a baseline comparison with a naive first-fit sequential placement of qubits relative to our structure-aware technique. We also do an architectural comparison between grids and linear architectures. We evaluate routing overhead using the inter-module instruction count, following \cite{sethi2026optimizing}. 
    
    \item We implement Toffoli gates via $\ket{\mathrm{CCZ}}$ injection, map the protocol to bicycle instructions including factory-to-module routing, estimate the resulting error with \texttt{bicycle\_numerics}, and model circuit failure using a Poisson distribution.
    \end{itemize}
    % \siyi{Very long very detailed, but lack of highlights, try to add \textbf{AAA} to highlight the real contribution.}


%
\section{Background}
\label{sub:back}
In this section, we proceed with the preliminaries.
% \siyi{Can devide it into section 2? how come intro so long.. :O}
\subsection{Multi-controlled Toffoli gates}
The doubly controlled X gate, commonly known as the Toffoli gate, is one of the most important reversible operators in quantum computing. Acting on three qubits, the basic Toffoli gate flips the target qubit if and only if both control qubits are in the state $\ket{1}$. Its generalization, the multi-controlled Toffoli ($n$-MCT) gate, consists of $n$ control qubits and one target qubit and performs a conditional bit flip on the target when all control qubits are simultaneously in the state $\ket{1}$. For more details on quantum gates, we refer the reader to~\cite{nc}.

Over the past few decades, numerous efforts have been made to reduce the resource requirements for implementing MCT gates. For basic Toffoli gates, state-of-the-art techniques employ measurement-based uncomputation to achieve decompositions using only four T gates~\cite{gidney2018quantum, jaques2020eurocrypt}. On the other hand, recent works on MCT decomposition employ the conditionally clean ancilla technique~\cite{nie2024arxiv, khattar2025quantum, dutta2025pra} to reduce the ancilla count and the circuit depth. In this technique, the availability of clean ancillae is assumed to be limited; therefore, some of the working qubits are made to be conditionally clean and treated as ancillae during the decomposition. The technique was first introduced by Nie et al.~\cite{nie2024arxiv}, formalized by Khattar et al.~\cite{khattar2025quantum}, and later generalized in Dutta et al.~\cite{dutta2025pra}. Moreover, Dutta et al.~\cite{dutta2025pra} also provided an $\lceil\log_2 n\rceil$ Toffoli-depth decomposition of an $n$-MCT using $n-2$ additional ancilla qubits, and proved its optimality by showing that the Toffoli depth cannot be reduced further, irrespective of the number of available ancilla.

The optimal-Toffoli-depth decompositions of $n$-MCT gates, assuming all-to-all qubit connectivity, for $n=5$ is shown in Figure~\ref{fig:optimal_mct}. For a detailed summary of state-of-the-art results on decomposing MCT gates into basic Toffoli gates, we refer the reader to \cite[TABLE II]{dutta2025pra}.
\begin{figure}[htbp]
    \centering
    \scalebox{0.9}{\includegraphics[scale=0.8]{Figures/n-mct.pdf}}
    % \scalebox{0.8}{\includegraphics[scale=0.7]{Figures/mct7.pdf}}
    \caption{Optimal Toffoli-depth $n$-MCT decompositions assuming all-to-all qubit connectivity, for $n=5$.}
    \label{fig:optimal_mct}
    % \vspace{-3.4em}
\end{figure}
% \begin{figure}[htbp]
% \centering
% \begin{subfigure}
%     \centering
%     \includegraphics[width=1\linewidth]{Figures/Logical_AND.pdf}
%     \caption{}
%     \label{fig:logicalAND}
% \end{subfigure}
% \begin{subfigure}
%     \centering
%     \includegraphics[width=1\linewidth]{Figures/mathias.pdf}
%     \caption{}
%     \label{fig:Jaques}
% \end{subfigure}
% \caption{Logical circuits for basic Toffoli decompositions with measurement-based uncomputation: (a) without ancilla \cite{gidney2018quantum}, and (b) with a single reusable ancilla \cite{jaques2020eurocrypt}.}
% \label{fig:toffoli}
% \end{figure}
%\anik{Need to give background of n-MCT decomposition directly.. can add the conditionally clean and the binary tree decomposition images}
%
% \paragraph{Fault-Tolerant Quantum Computation}
% Quantum information is inherently susceptible to decoherence and operational errors. Fault-tolerant quantum computation addresses these challenges by encoding logical qubits into multiple physical qubits using quantum error-correcting codes. Periodic syndrome measurements identify error syndromes without disturbing the encoded quantum information, enabling reliable quantum computation even in the presence of physical noise.

% The execution of an encoded quantum circuit is considerably more involved than its logical gate sequence. Each logical operation requires syndrome extraction, decoder processing, and error recovery, while non-Clifford gates frequently rely on auxiliary procedures such as magic-state preparation and distillation. Consequently, the physical cost of a logical circuit is determined by multiple interacting factors, including encoding overhead, stabilizer measurement complexity, decoder latency, and logical gate implementation.

% Accordingly, evaluating a quantum circuit solely through logical metrics such as T-count, T-depth or ancilla count provides an incomplete picture of its practical execution cost. A comprehensive resource estimation framework must incorporate both logical circuit complexity and fault-tolerant implementation overhead.

% %
% \paragraph{qLDPC Codes}
% Quantum low-density parity-check (qLDPC) codes constitute an important family of stabilizer codes characterized by sparse parity-check matrices. Similar to their classical counterparts, each stabilizer generator acts on only a constant number of qubits, and each qubit participates in a bounded number of stabilizer measurements. This sparsity enables efficient syndrome extraction and offers the potential for improved scalability compared with conventional topological codes.

% Recent advances have produced several families of qLDPC codes with asymptotically favorable parameters, including hypergraph-product codes, balanced-product codes, lifted-product codes, fiber-bundle codes, and quantum Tanner codes. These constructions demonstrate that constant-rate quantum codes with linear minimum distance are achievable, representing a major milestone in quantum coding theory.

% Although surface codes offer a high error threshold, qLDPC codes promise substantially lower physical-qubit overhead while preserving fault tolerance. However, their practical performance depends on efficient decoding, logical gate implementation, and communication strategies, all of which significantly influence the resource requirements of encoded quantum circuits.

% %
% \paragraph{Resource Estimation for Fault-Tolerant Circuits}
% Beyond logical gate counts, practical resource estimation must account for physical qubit overhead, code distance, syndrome extraction cycles, decoder latency, logical execution time, and space-time volume. For MCT gates, improvements in logical metrics do not necessarily translate into lower physical resource requirements. For instance, two decompositions with identical T-counts may incur significantly different costs due to variations in syndrome extraction, encoded communication, and logical gate scheduling under a given quantum error-correcting code.

% Consequently, architecture- and code-aware resource estimation is essential for evaluating the practicality of large-scale quantum computation. In this paper, we develop a qLDPC-aware analytical framework that maps state-of-the-art logical MCT decompositions to their corresponding fault-tolerant resource requirements, enabling a quantitative comparison of different decomposition strategies.
%
%\cite{gidney2021cccz}
\subsection{qLDPC and BB codes}
Quantum low-density parity-check (qLDPC) codes constitute an important family of stabilizer codes characterized by sparse parity-check matrices. Similar to their classical counterparts, qLDPC codes have the property that each stabilizer generator acts on only a constant number of qubits, while each qubit participates in a bounded number of stabilizer checks. This sparsity enables efficient syndrome extraction and, more importantly, allows qLDPC codes to achieve favorable encoding rates and code distances with substantially lower physical-qubit overhead than conventional topological codes. Recent constructions, including hypergraph-product codes~\cite{tillich2013tit}, quantum Tanner codes~\cite{leverrier2022focs}, balanced-product codes~\cite{bpqc2021}, and lifted-product codes~\cite{lpqc}, have demonstrated the possibility of constructing quantum codes with constant encoding rate and asymptotically growing distance.
\anik{I am thinking of removing the examples due to space constraints}
% Figure~\ref{fig:qldpc-flowchart} summarizes different types of qLDPC codes, with representative examples. 
% \begin{figure*}[ht]
%     \centering
%     \scalebox{0.95}{\input{Figures/flowchart}}
%     \caption{Different varieties of existing qLDPC codes \siyi{limited info figure, ummmm try to improve or delete}}
%     \label{fig:qldpc-flowchart}
% \end{figure*}

% Unlike topological codes such as the surface code~\cite{surface-code}, qLDPC codes are not necessarily defined by geometrically local interactions. 
The sparse connectivity specified by their parity-check matrices can instead induce long-range interactions between physical qubits. Consequently, implementing qLDPC codes requires not only efficient syndrome extraction and decoding but also hardware architectures \cite{wang2025nature} and compilation strategies \cite{tourdegross}, similar to lattice surgery in surface codes ~\cite{gameofsurfacecodes,leblond2024realistic} that support their non-local connectivity. Among the various qLDPC constructions, bivariate bicycle (BB) codes are particularly relevant to practical fault-tolerant quantum computing. Recent experimental demonstrations have shown that BB codes can be implemented on superconducting quantum hardware by exploiting hardware capabilities that support long-range interactions~\cite{bravyi2024highthreshold}.

% BB codes are CSS qLDPC codes constructed from two sparse bivariate polynomials over $\mathbb{F}_2$. Let $a(x,y)$ and $b(x,y)$ be two such polynomials, and let $S_\ell$ and $S_m$ denote cyclic-shift matrices. Evaluating the polynomials at these shift matrices gives two binary matrices 
% $A = a(S_\ell, S_m),\quad B = b(S_\ell, S_m).$
% Here, the variables $x$ and $y$ correspond to cyclic shifts of the physical qubits along two dimensions. Since the corresponding cyclic-shift matrices commute, the matrices $A$ and $B$ also commute. The CSS parity-check matrices are then defined as
% $$H_X = [A \; B], \quad H_Z = [B^T \; A^T].$$
% The stabilizer commutation condition directly follows from
% $$H_XH_Z^T=[A \; B]
% \begin{bmatrix}B \\A \end{bmatrix} = AB+BA = 0 \pmod 2,$$
% where the last equality follows from the commutativity of $A$ and $B$ over $\mathbb{F}_2$. Thus, the rows of $H_X$ and $H_Z$ define mutually commuting $X$- and $Z$-type stabilizer generators.

% The sparsity of the polynomials $a(x,y)$ and $b(x,y)$ is inherited by $A$ and $B$, and consequently by $H_X$ and $H_Z$. Therefore, each stabilizer generator acts on only a bounded number of physical qubits, while each physical qubit participates in a bounded number of stabilizer checks, establishing the LDPC property of the code. Importantly, these sparse connections need not be geometrically local in a two-dimensional nearest-neighbor layout. Instead, the underlying cyclic-shift structure naturally induces long-range connections in the physical Tanner graph. It also provides a set of code-preserving automorphisms, allowing physical and logical qubit labels to be permuted according to symmetries of the code without changing its stabilizer structure. These properties make BB codes particularly interesting for fault-tolerant architectures with long-range connectivity. For a CSS code with $n$ physical data qubits, the number of encoded logical qubits is $k = n-\operatorname{rank}(H_X)-\operatorname{rank}(H_Z).$ In contrast to a standard surface-code patch, which typically encodes a single logical qubit using $O(d^2)$ physical qubits, a BB-code block can encode multiple logical qubits within a single block. For example, the $\left[\left[144,12,12\right]\right]$ BB/Gross-code instance encodes $12$ logical qubits into $144$ physical data qubits and has code distance $12$. One of the 12 logical qubits is reserved as a pivot/ancilla qubit for routing and logical operations. The high encoding rate, together with the structured long-range connectivity of BB codes, makes them a promising candidate for fault-tolerant implementations of large quantum circuits.
BB codes are CSS qLDPC codes constructed from two sparse bivariate polynomials $a(x,y)$ and $b(x,y)$ over $\mathbb{F}_2$. Evaluating these polynomials on cyclic-shift matrices $S_\ell, S_m$ provides two binary matrices $A = a(S_\ell,S_m),\, B = b(S_\ell, S_m)$, from which the CSS parity check matrices are defined as
$$H_X = [A \; B], \,H_Z = [B^T \; A^T].$$
The sparsity of $a$ and $b$ yields sparser stabilizer checks, thereby preserving the LDPC property. Unlike surface-code stabilizers, these interactions need not be geometrically local in a 2D nearest-neighbor layout. The cyclic-shift structure naturally induces structured long-range connectivity. BB codes also encode multiple logical qubits within a single code block, providing substantially higher encoding rates than conventional surface-code patches. In this work, we consider the $\left[[144,12,12]\right]$ BB-code instance, which encodes 12 logical qubits into 144 physical data qubits with distance 12. One logical qubit is reserved as a pivot/ancilla for logical operations and routing.

%%%%%%%%%%%%%%%%%%%%%%%%%%%%%%
% As mentioned in Section \ref{sec:intro}, qLDPC codes offer improvement to surface codes because of their encoding rate. Like the hamming code \cite{}, any quantum error correcting code can be represented by a Tanner graph, the only difference in quantum codes is that the checks are for errors on the 2 pauli bases $X,Z$. This checking overhead typically makes quantum codes resource-heavy and harder to execute; however this also brings forth a new paradigm for hardware and raises new compilation concerns \cite{gameofsurfacecodes,tourdegross,gidney2025factor,leblond2024realistic}

% There exists a variety of QLDPC codes in theory like the hypergraph product code \cite{tillich2013tit}, quantum tanner codes \cite{leverrier2022focs}. But the BB Code has been proved to be executed on the IBM superconducting hardware \cite{bravyi2024highthreshold} due to their new hardware capabilities for long range connections. Bivariate bicycle (BB) codes are CSS quantum LDPC codes defined from two bivariate polynomials over \(\mathbb{F}_2\).  Let \(a(x,y)\) and \(b(x,y)\) be sparse polynomials, and let \(S_\ell\) and \(S_m\) denote cyclic shift matrices. The
% polynomials define two binary matrices \(A = a(S_\ell,S_m), \qquad B = b(S_\ell,S_m),\). In the polynomials $x,y$ correspond to a cyclic shift of the data qubits. Since these are commuting in nature the matrices \(A,B\) also commute. The CSS parity check matrices are now 
% \[
% H_X = [A \; B], \qquad H_Z = [B^T \; A^T].
% \]
% The stabilizer commutation condition follows as 
% \[
% H_XH_Z^T
% =
% [A \; B]
% \begin{bmatrix}
% B \\
% A
% \end{bmatrix}
% =
% AB+BA
% =
% 0 \pmod 2.
% \]
% Thus, the rows of \(H_X\) and \(H_Z\) define mutually commuting \(X\)- and \(Z\)-type stabilizer generators. 

% The sparsity of \(A,B\) also reflects on \(H_X, H_Z\), so each stabiliser check only acts on a bounded number of physical qubits, this reflects the LDPC property of the code. These sparse connections are not necessarily geometrically local in a
% simple two-dimensional nearest-neighbor layout. In stead the cyclic shift structure allows long-range physical Tanner-graph connections and automorphisms, i.e physical/logical qubit labels can be permuted according the symmetries of the code preserving the code structure. 

% A key advantage of BB codes over the surface code is their encoding efficiency.
% For a CSS code with \(n\) physical data qubits, the number of encoded logical
% qubits is
% \[
% k = n-\operatorname{rank}(H_X)-\operatorname{rank}(H_Z).
% \]
% While a standard surface-code patch typically encodes one logical qubit using
% \(O(d^2)\) physical qubits, a BB-code block can encode multiple logical qubits
% in a single block. For example, the \([[144,12,12]]\) BB/Gross-code instance
% encodes \(12\) logical qubits into \(144\) physical data qubits with distance
% \(12\). 
 
% \anik{Basic math of BB code}
\subsection{Fault Tolerant compilation for BB codes}
Compilation for fault-tolerant architectures differs fundamentally from compilation for NISQ-based architectures~\cite{bhaumik2026optimal}. Rather than targeting individual physical qubits, fault-tolerant compilation operates at the logical level, where physical qubits are organized into encoded logical blocks. In the surface code, a logical qubit is typically represented by a code patch, whereas in BB-code-based architectures, a single code block or module can encode multiple logical qubits. Consequently, the target instruction set may also differ substantially from a conventional gate-based instruction set. A compiler may instead map the algorithm onto the fault-tolerant logical operations supported by the underlying architecture.

Another important aspect of fault-tolerant compilation is the placement and routing of magic-state factories~\cite{gameofsurfacecodes}. Non-Clifford operations, such as T gates, typically require the preparation and injection of magic states to implement the corresponding logical operation. Magic-state factories continuously produce these resource states, introducing an additional architectural component that must be considered during compilation. In particular, the compiler must account for the placement of factories and the communication required to deliver magic states to the logical qubits that consume them. In this paper, we consider a $\ket{\mathrm{CCZ}}$-state injection protocol for executing our Toffoli-intensive workloads. This approach allows the non-Clifford resource generation to be largely separated from the computation itself and enables optimizations at the logical level of the Toffoli network.

Yoder et al.~\cite{tourdegross} recently introduced a compilation scheme for BB-code-based quantum hardware. They developed a new instruction set architecture (ISA) for compiling quantum instructions to a fault-tolerant abstraction, together with corresponding logical error-rate estimates. These compilers currently support results through simulation as cloud-based FTQC hardware are not available yet. 

% \subsection{Organization and Contributions}
% \label{sub:org}
% In this section, we summarize section-wise contributions along with an overall structure of this paper.



%%% NEW SECTION
% \section {}
\section{Tree based Placement in BB Code Architecture}
\label{sec:1}
% \siyi{QUestion: should this section just be the structure of BB with a focus on the Placement}
In this work, we abstract each decoded BB-code block as a logical module with capacity $C=11$ and study how Toffoli/CCZ interactions should be placed and routed across a graph of such modules. The BB-code-based architectures are not yet well established. Yoder et al.~\cite[Figure~2(a)]{tourdegross} considered a linear architecture in which BB modules are arranged sequentially, with a magic-state factory placed at one end. Sethi et al.~\cite{sethi2026optimizing} extended this architecture to a long-grid topology.

From~\cite[Table~2]{tourdegross}, it is evident that the cost of inter-module operations is substantially higher than that of intra-module operations. Hence, we use the number of inter-module interactions as our primary routing metric. The optimal Toffoli-depth decomposition of an $n$-MCT~\cite{dutta2025pra} has a hierarchical structure composed of smaller $k$-MCTs. We exploit this structure through a tree-based placement strategy that groups the $k$-MCTs corresponding to the same lower-level subtree into a single BB module. The intermediate ancilla qubits created by each group are then propagated toward the higher-level outputs. This preserves the hierarchical dependencies of the decomposition while reducing unnecessary inter-module communication. The grouping factor $k$ determines both the number of controls assigned to each module and the number of BB modules required, given by $\lceil  n/k \rceil$. Consequently, the choice of $k$ introduces a topology-dependent trade-off among local packing efficiency, communication distance, and residual ancilla capacity.

Each module has a maximum capacity of 11 logical qubits and can therefore accommodate at most $k=6$ control qubits (as the optimal Toffoli-depth $n$-MCT decomposition requires $2n-1$ working qubits~\cite{dutta2025pra}). However, accounting for the additional ancilla required to compose the complete MCT may make such a choice impractical. In particular, it may require additional uncomputation, increasing the gate complexity, or leave insufficient capacity for dirty ancillae. We therefore derive a bound on the permissible value of $k$ in Proposition~\ref{prop:prop1}.

\begin{proposition}
\label{prop:prop1}
   Consider an $n$-MCT decomposition mapped onto BB-code-based modules with logical capacity $C=11$. Then, each module can implement a maximum of $k$-MCT for $k \leq 5$, where $\lceil  n/k \rceil$ modules provide sufficient residual logical-qubit capacity to accommodate the higher-level MCT operations in the decomposition without requiring additional ancilla qubits.
\end{proposition}
\begin{proof}
From~\cite{dutta2025pra}, the binary-tree-based decomposition of an $n$-MCT requires $2n-1$ working qubits in total, including $n-2$ ancilla. Since each BB-code-based module has logical capacity $C=11$, a single module can theoretically accommodate a $6$-MCT, which requires 11 qubits. However, using the full module capacity for the lower-level MCTs leaves insufficient capacity for the ancillae required to construct larger MCTs. To avoid additional inter-module operations and ancilla management, we restrict the maximum MCT size per module to $n=5$. Hence, $k\leq 5$.

Consequently, implementing an $n$-MCT with $n\gg 5$ requires at least $\lceil n/5 \rceil$ many BB-code-based modules to accommodate the decomposition without requiring additional ancilla qubits.
\end{proof}

We find that both the number and placement of magic-state factories significantly influence routing efficiency, as each gate requiring a magic-state resource must communicate with a factory. Consequently, the factory configuration affects routing at a finer-grained, gate-level scale.

%We find that both the number and placement of factories significantly influence routing efficiency, since each gate requiring a magic-state resource must communicate with a factory. Consequently, factory configuration affects routing at a finer-grained, gate-level scale.

% \paragraph{Calculation of Routing Volume}

% \subsection{Tree-based placement}
% \anik{Need a diagram depicting tree-based implementation of a binary tree and its mapping on the BB module}

\subsection{Linear Topology placement} 
\label{seb:lin-top}
We first consider a linear architecture in which the BB modules are arranged sequentially along a single line. This is one of the most commonly explored architectures. Magic-state factories are attached to the modules at regular intervals, and we vary the factory placement period to study the effects of factory density and placement on routing performance. Figure~\ref{fig:routing_volume_comparison} shows how different $k$-groupings affect the number of inter-module interactions under a first-fit-based placement strategy. Varying the factory-placement period also narrows the gap between the different groupings and the baseline.
\begin{figure}[htbp]
    \centering
\scalebox{0.7}{    \includegraphics[width=0.9\linewidth]{Figures/linvsgrid-1.pdf}
}    \caption{Linear vs. grid factory placement, showing heuristic-selected factory modules in red and standard modules in green.}
    \label{fig:linvsgrid}
\end{figure}

\begin{figure*}[htbp]
    \centering

    % ============================================================
    % Linear architecture: varying factory period
    % ============================================================
    
    \begin{subfigure}[t]{0.32\textwidth}
        \centering
        \includegraphics[width=\linewidth]
        {Figures/routing_volume_comparison/routing_volume_comparison_period_2.pdf}
        \caption{$\mathrm{Period}=2$}
        \label{fig:routing_period_2}
    \end{subfigure}
    \hfill
    \begin{subfigure}[t]{0.32\textwidth}
        \centering
        \includegraphics[width=\linewidth]
        {Figures/routing_volume_comparison/routing_volume_comparison_period_3.pdf}
        \caption{$\mathrm{Period}=3$}
        \label{fig:routing_period_3}
    \end{subfigure}
    \hfill
    \begin{subfigure}[t]{0.32\textwidth}
        \centering
        \includegraphics[width=\linewidth]
        {Figures/routing_volume_comparison/routing_volume_comparison_period_4.pdf}
        \caption{$\mathrm{Period}=4$}
        \label{fig:routing_period_4}
    \end{subfigure}

    \vspace{0.6em}

    % ============================================================
    % Grid architecture: varying grouping size k
    % ============================================================

    \begin{subfigure}[t]{0.32\textwidth}
        \centering
        \includegraphics[width=\linewidth]
        {Figures/routing_volume_grid_improvement/k3.pdf}
        \caption{$k=3$}
        \label{fig:grid_routing_k3}
    \end{subfigure}
    \hfill
    \begin{subfigure}[t]{0.32\textwidth}
        \centering
        \includegraphics[width=\linewidth]
        {Figures/routing_volume_grid_improvement/k4.pdf}
        \caption{$k=4$}
        \label{fig:grid_routing_k4}
    \end{subfigure}
    \hfill
    \begin{subfigure}[t]{0.32\textwidth}
        \centering
        \includegraphics[width=\linewidth]
        {Figures/routing_volume_grid_improvement/k5.pdf}
        \caption{$k=5$}
        \label{fig:grid_routing_k5}
    \end{subfigure}


    \begin{subfigure}[t]{0.32\textwidth}
            \centering
            \includegraphics[width=\linewidth]
            {Figures/error_estimations/results_gross_1e-3_4_total_error.pdf}
            \caption{Aggregate logical error at $p=10^{-3}$.}
            \label{fig:gross_1e3_aggregate_error}
        \end{subfigure}
        \hfill
        \begin{subfigure}[t]{0.32\textwidth}
            \centering
            \includegraphics[width=\linewidth]
            {Figures/error_estimations/results_gross_1e-3_4_failure_probability.pdf}
            \caption{Circuit failure probability at $p=10^{-3}$.}
            \label{fig:gross_1e3_failure_probability}
        \end{subfigure}
        \hfill
        \begin{subfigure}[t]{0.32\textwidth}
            \centering
            \includegraphics[width=\linewidth]
            {Figures/error_estimations/results_gross_1e-4_4_failure_probability.pdf}
            \caption{Circuit failure probability at $p=10^{-4}$.}
            \label{fig:gross_1e4_failure_probability}
        \end{subfigure}
    \caption{
        Inter-module routing characteristics under different factory
        placements and grouping configurations. 
        (a)--(c) show the linear BB-module architecture for factory
        periods $2$, $3$, and $4$, respectively, while
        (d)--(f) show the improvement in grid architecture vs linear for grouping factors $k=3$, $4$, and $5$.
        (g)--(h) shows the aggregate and circuit failure probabilities for hardware error rates of $p = 10^{-3}$ and and $p = 10 ^{-4}$
    }
    \label{fig:routing_volume_comparison}
\end{figure*}


\subsection{Grid Topology Placement}
\label{sub:grid-top}
We find that in the grid-based structure, factory density and factory placement take precedence on improving routing overhead. Sethi et al.~\cite{sethi2026optimizing} previously studied BB-code placement on grid topologies. However, their architecture places all the magic-state factories along a single edge of the grid, which does not adequately capture the distributed factory configurations considered in our work.

To investigate the impact of factory placement, we consider a square-grid architecture with dimensions $p = \sqrt{M}$ and $q = \lceil M/p \rceil$, where $M$ denotes the total number of BB modules. We compare the grid and linear topologies while keeping the number of magic-state factories fixed. For the grid topology, we use a facility-location heuristic to optimize factory placement and evaluate its impact on Toffoli routing performance, an example of such a placement is shown in Figure ~\ref{fig:linvsgrid}

\subsubsection{Facility Location Optimization Heuristic}
\label{subsub:floh}
This factory placement problem resembles facility location. We initialize factory sites using graph medians and refine them with swap-based local search. Each factory is placed in a separate module and connected to its assigned BB module by an additional edge, reflecting their treatment as distinct entities in fault-tolerant compilation. To determine the factory locations, we first employ a greedy heuristic~\cite{hakimi1964optimum}, followed by a 1-swap local-search heuristic~\cite{teitz1968heuristic}.

Let $G=(V, E)$ denote the BB-module graph, and let $d(u,v)$ denote the shortest-path distance between vertices $u,v\in V$. Given a set $S\subseteq V$ of factory sites, the greedy placement strategy minimizes the objective
$$J(S) = \sum_{v \in V} \min_{s\in S}d(v,s); \quad S \subseteq V,$$
where $S$ denotes the subset of modules selected as factory sites. After the greedy procedure has considered all vertices according to this objective, we perform a 1-swap local search over the candidate set. Specifically, each selected factory site $s \in S$ is tentatively replaced by an unselected module $v \in V \setminus S$. The replacement is accepted whenever $J\!\left((S \setminus \{s\}) \cup \{v\}\right) < J(S).$
The procedure terminates when no improving single-site swap exists.

% For the logical Clifford+T decomposition of basic Toffoli gate, suppose the resource metrics are denoted by logical qubit ($Q_L$), T-count ($\mathrm{T}_c$), T-depth ($\mathrm{T}_d$), Clifford gates (C). Further, suppose the BB code requires $n$ many physical qubit ($Q_P$) to correct one single logical qubit, thus, $Q_P = nQ_L$. 

%%% NEW SECTION
\subsection{Evaluation of Routing Results}
\label{sub:evaluation}
% \anik{need to mention that maximum fit first placement is done}
% \anik{shorten this}
We model the architectures as graphs and simulate qubit placement and instruction flow to quantify inter-module routing count and distance. The routing metric quantifies the number of inter-module interactions present, as previously modeled by \cite{sethi2026optimizing}. Specifically, for each interaction, we compute the distance between the factory and the modules containing the corresponding data qubits and sum these distances over all interactions. For example, consider a Toffoli gate $(c_1,c_2,t)$ and a $\ket{\mathrm{CCZ}}$ factory $F$. Since consuming a $\ket{\mathrm{CCZ}}$ state requires three independent Pauli-product measurements for injection. The number of inter-module instructions is given by
$$R = d(M(c_1),F) + d(M(c_2),F) + d(M(t),F)$$
where $M(q)$ denotes the module containing the qubit $q$. 

We compare our placement strategy against a first-fit baseline, in which the qubits are placed sequentially as shown in Figure~\ref{fig:optimal_mct}, following the order $c_1,c_2,\ldots , c_n,a_1,a_2,\ldots ,a_{n-2},t$. Then we compare grouping $k$-MCTs. Specifically, we partition $n$ control qubits into $\lceil n/k \rceil$ modules, form a $k$-MCT within each module, and ensure that the remaining ancillae can be accommodated within the residual capacity of each module in a first-fit fashion. We evaluate both the placement strategies for $n$ up to 127.

%
\subsubsection{Grouping and Factory-Period in Linear Topology}
Figures~\ref{fig:routing_period_2}, \ref{fig:routing_period_3}, and \ref{fig:routing_period_4} show the routing performance for different factory periods, where the factory period determines the spacing between consecutive factories along the linear topology. Across all three cases, the tree-based grouping strategy consistently improves the routing metric relative to the baseline. As the factory period increases, however, the gap between the baseline and the grouping strategy decreases, indicating that the efficiency of the linear topology deteriorates with increasing factory spacing. The maximum mean improvement is by $16.02\%$ for, $period = 2$. 

% \begin{table}[h]
%     \centering
%     \begin{tabular}{|c|c|c|c|}
%         \hline
%         $Period$ & $k = 3$ & $k = 4$ & $k = 5$ \\
%         \hline
%         2 & \cellcolor{green!35}10.25\% 
%           & \cellcolor{green!55}14.14\% 
%           & \cellcolor{green!70}\textbf{16.02\%} \\
%         \hline
%         3 & \cellcolor{yellow!35}6.72\% 
%           & \cellcolor{green!40}11.45\% 
%           & \cellcolor{green!50}13.22\% \\
%         \hline
%         4 & \cellcolor{red!20}4.64\% 
%           & \cellcolor{yellow!45}8.77\% 
%           & \cellcolor{green!35}10.26\% \\
%         \hline
%     \end{tabular}
%     \caption{Mean improvement over 128-control MCT mapped in linear BB-module architecture over $k$.}
%     \label{tab:linear_improvement_k}
% \end{table}
%
\subsubsection{Comparison of Grid Topology vs Linear Topology}
We compare the linear BB-module topology with a grid topology that has the same number of modules, $M$, and the same number of magic-state factories, as determined by the factory period, $f$. Figures~\ref{fig:grid_routing_k3}, \ref{fig:grid_routing_k4}, and \ref{fig:grid_routing_k5} show that relative improvement decreases modestly with k, from a maximum of $27.7\%$ at $k=3$ to $23.7\%$ at $k=5$. Tree-based placement has transferred inefficiently, since it assumes subtrees occupy the contiguous maximum first fit, which terms as locality on a chain but not on a grid. But, since $k$ is bounded in Proposition~\ref {prop:prop1}, these values cover the full set of admissible grouping factors ($k$), and the grid topology provides a positive impact throughout this range of $k = [3,5]$.

However, each of the figures also show that the improvement in the routing metric provided by the grid topology increases monotonically with $f$, ranging from $6.1\%$ to $27.7\%$ for $k=3$. In other words, the grid's advantage grows as factories thin out, this is where the linear's diameter starts to dominate. 
 
% \vspace{-2em}
%
\section{Error Estimation for Routing}
\label{sec:2}
%Methodology
Error is the primary attribute that FTQC architectures claim to optimize, Yoder et al.~\cite[Table 2]{tourdegross} provides a bicycle ISA for compiling quantum circuits onto BB-code-based architectures. They also characterize the error rates associated with each of these operations. Hence, we develop a workflow that takes the architecture details of the mapped Toffoli circuit and its gate-DAG as input, generates a sequence of Bicycle-ISA instructions, and feeds the resulting instructions into the error estimator \texttt{bicycle\_numerics}. The Toffoli benchmarks use $\ket{\mathrm{CCZ}}$ injection  which is not natively supported by \texttt{bicycle\_compiler}~\cite{bicycle_compiler}.
. We therefore model the instructions required for the injection's communication and correction rounds explicitly for each Toffoli, and accumulate the resulting errors.
% \begin{figure}[htbp]
%     \centering
% \scalebox{0.8}{    \includegraphics[width=\linewidth]{Figures/ccz_injection.pdf}
% }    \caption{CCZ injection protocol}
%     \label{fig:cczinjection}
% \end{figure}

% \begin{figure*}[htpb]
%     \centering

%     % =====================================
%     % p = 10^{-3}
%     % =====================================
%     \begin{subfigure}[t]{0.46\textwidth}
%         \centering
%         \scalebox{0.85}{\includegraphics[width=0.88\linewidth]
%         {Figures/error_estimations/results_gross_1e-3_4_total_error.pdf}}
%         \caption{Aggregate logical error at $p=10^{-3}$.}
%         \label{fig:gross_1e3_aggregate_error}
%     \end{subfigure}
%     % \hfill
%     \begin{subfigure}[t]{0.46\textwidth}
%         \centering
%         \scalebox{0.85}{\includegraphics[width=0.88\linewidth]
%         {Figures/error_estimations/results_gross_1e-3_4_failure_probability.pdf}}
%         \caption{Estimated circuit failure probability at $p=10^{-3}$.}
%         \label{fig:gross_1e3_failure_probability}
%     \end{subfigure}

%     % \vspace{0.4em}

%     % =====================================
%     % p = 10^{-4}
%     % =====================================
%     \begin{subfigure}[t]{0.46\textwidth}
%         \centering
%         \scalebox{0.85}{\includegraphics[width=0.88\linewidth]
%         {Figures/error_estimations/results_gross_1e-4_4_total_error.pdf}}
%         \caption{Aggregate logical error at $p=10^{-4}$.}
%         \label{fig:gross_1e4_aggregate_error}
%     \end{subfigure}
%     % \hfill
%     \begin{subfigure}[t]{0.46\textwidth}
%         \centering
%         \scalebox{0.8}{\includegraphics[width=0.88\linewidth]
%         {Figures/error_estimations/results_gross_1e-4_4_failure_probability.pdf}}
%         \caption{Estimated circuit failure probability at $p=10^{-4}$.}
%         \label{fig:gross_1e4_failure_probability}
%     \end{subfigure}

%     \caption{
%         Aggregate logical error and circuit failure probability for
%         decomposed-MCT circuits mapped to the Gross-code architecture
%         at physical error rates $p=10^{-3}$ and $p=10^{-4}$.
%         Figures~(a) and (c) show the aggregate logical error,
%         $E=\sum_i p_i$, while Figures~(b) and (d) show the corresponding
%         failure probability,
%         $P_{\mathrm{fail}} = 1-e^{-E}$.
%     }
%     \label{fig:gross_error_characteristics}
% \end{figure*}
% \subsection{Error Calculation}
% \label{sub:error_calc}
% The CCZ-injection protocol is not modeled directly by the estimator, which instead assumes T-state injection. Since our benchmarks consist entirely of Toffoli gates, CCZ-state injection is a critical component of the overall error model. We therefore divide the error associated with the CCZ-injection process into three components: factory cost, injection cost, and correction cost, as illustrated in Figure~\ref{fig:cczinjection}.

%
% \begin{figure*}[htpb]
%     \centering

%     \begin{subfigure}[t]{0.32\textwidth}
%         \centering
%         \includegraphics[width=\linewidth]
%         {Figures/error_estimations/results_gross_1e-3_4_total_error.pdf}
%         \caption{Aggregate logical error at $p=10^{-3}$.}
%         \label{fig:gross_1e3_aggregate_error}
%     \end{subfigure}
%     \hfill
%     \begin{subfigure}[t]{0.32\textwidth}
%         \centering
%         \includegraphics[width=\linewidth]
%         {Figures/error_estimations/results_gross_1e-3_4_failure_probability.pdf}
%         \caption{Circuit failure probability at $p=10^{-3}$.}
%         \label{fig:gross_1e3_failure_probability}
%     \end{subfigure}
%     \hfill
%     \begin{subfigure}[t]{0.32\textwidth}
%         \centering
%         \includegraphics[width=\linewidth]
%         {Figures/error_estimations/results_gross_1e-4_4_failure_probability.pdf}
%         \caption{Circuit failure probability at $p=10^{-4}$.}
%         \label{fig:gross_1e4_failure_probability}
%     \end{subfigure}

%     \caption{
%         Aggregate logical routing-error and circuit failure probability for
%         decomposed-MCT circuits mapped to the Gross-code architecture.
%         (a) shows the aggregate logical error at $p=10^{-3}$, while
%         (b) and (c) show the circuit routing failure probabilities at
%         $p=10^{-3}$ and $p=10^{-4}$, respectively. \siyi{Nice figure! but difficult to distinguish the dash grey columns}
%     }
%     \label{fig:gross_error_characteristics}
% \end{figure*}
\paragraph{Factory distillation cost}
\label{subsub:fact_dist_cost}
The factory distillation cost corresponds to the cost of generating a $\ket{\mathrm{CCZ}}$ state in a magic-state factory. We assume this cost to be constant for a given factory, with the corresponding cost incurred each time a $\ket{\mathrm{CCZ}}$ state is produced. In this study, however, we omit the distillation cost and focus primarily on errors arising from routing and inter-module communication.

%
\paragraph{Injection cost}
\label{subsub:inj_cost}
The injection cost arises from entangling the $\ket{\mathrm{CCZ}}$ state with the data qubits and subsequently measuring the corresponding factory qubits. We model this process using \texttt{Measure} operations within the modules and \texttt{JointMeasure} operations across the logical qubits of different modules as shown in \cite{bicycle_compiler}

At the gate level, the injection procedure contains a $CNOT_{d_i\rightarrow r_i}$ followed by a measurement of the $r_i$ qubit, which is equivalent to a destructive two-qubit Pauli-product measurement, $Z_{d_i}Z_{r_i}$. For each Toffoli gate, we first select the nearest available factory by evaluating the distances between the factory and the modules containing the participating data qubits. After selecting the factory, we decompose the corresponding injection operations. For each $Z_{d_i}Z_{r_i}$ measurement, we consider two cases:
\begin{itemize}
	\item If $d_i$ is located in the same module to which the factory is attached, we model the operation as a local \texttt{Measure} operation within that module.
	\item If $d_i$ is located in a different module from the factory, we perform a local \texttt{Measure} operation in each module along the communication path between the factory module and the module containing $d_i$, and connect these operations using a sequence of \texttt{JointMeasure} instructions.
    % \item $d_i$ is far away from the factory module
\end{itemize}
% \vspace{-1em}
%
\paragraph{Correction cost}
\label{subsub:corr_cost}
Once the entanglement and measurement steps are complete, conditional Clifford corrections are required to obtain the desired output state up to a global phase. The correction cost depends on the measurement outcomes $m_0, m_1, m_2$ of the factory qubits. As discussed in~\cite[Section~3.4]{tourdegross}, Clifford corrections that cross module boundaries can only be implemented through cross-module measurements. The conditional corrections in the CCZ-injection protocol consist of $CZ$ and $Z$ operations. Since a $CZ$ correction may require an inter-module measurement, we estimate the expected correction cost by considering the probability that each conditional $CZ$ correction is required.

Let $C_0$, $C_1$, and $C_2$ denote the costs associated with the three possible $CZ$ corrections. Assuming that each measurement outcome $m_i$ is independently equal to 1 with probability $1/2$, the expected correction cost is .
$$\mathbb{E}\!\left[C_{corr}\right] = \mathbb{E}\!\left[m_0 C_0 + m_1 C_1 + m_2 C_2\right] = \frac{1}{2}C_0 + \frac{1}{2}C_1 + \frac{1}{2}C_2.$$

For each inter-module $CZ$ correction, we generate the corresponding \texttt{JointMeasure} operation with a weight of $1/2$ in the expected error calculation. Single-qubit $Z$ corrections are omitted because their cost within a module is negligible compared with that of inter-module operations and therefore contributes negligibly to the overall error estimate.

%
\subsection{Error Evaluation}
\label{sub:error_evaluation}
% \anik{Indicate how ccz error injection has been estimated, injection cost is there but factory cost is not there which needs to be added at last}


We report two error metrics. The first is the aggregate error, \[E= \sum_i p_i\]
Where $p_i$ denotes the logical error probability associated with the $i$-th instruction. The quantity $E$ represents the expected number of logical faults per Toffoli-circuit execution. Because it is an additive aggregate of individual instruction error probabilities, $E$ grows approximately linearly with the number of error-prone operations. Importantly, $E$ is an expected fault count rather than a probability and may therefore exceed unity. Next, we estimate the circuit failure probability as 
\[P_{fail} = 1 - e^{-E}\]
This expression follows from modeling the number of faults as a Poisson random variable with mean $E$, under the assumption that instruction failures are approximately independent and that the individual error probabilities are sufficiently small. 
% The probability of observing at least one fault is then $1-e^{-E}$. Distillation error is excluded from this estimate because it is independent of qubit placement. 
Consequently, the reported figures quantify placement-dependent error rather than the end-to-end failure probability. We evaluate on the estimator provided by \cite{tourdegross,bicycle_compiler} which have encoded error rates for given instructions in the bicycle-ISA. The code for instruction generation is to be provided in GitHub \footnote{\url{https://github.com/anik314159/QLDPC-Toffoli}}.

\subsubsection{Results}
\label{subsub:result}

The error evaluations given are only for linear BB-code-based architectures. Although the Bicycle architecture permits more general module-based connectivity, extending the error analysis to the proposed two-dimensional grid would require additional assumptions regarding inter-module routing and physical connectivity. No significant improvement is seen over period, However in the aggregate error estimation (Figure \ref{fig:gross_1e3_aggregate_error}) the error-relations display decrease of error over increase of $k$ hence we chose an average value of $\textit{period}=4$, we report the aggregate error in Figure~\ref{fig:gross_1e3_aggregate_error}, and circuit failure probability for $p = 10^{-3}$ and $p = 10 ^{-4}$ Figure~\ref{fig:gross_1e3_failure_probability}, Figure~\ref{fig:gross_1e4_failure_probability}. Following ~\cite[Section~4]{tourdegross}, we use a circuit failure threshold of $1/3$. However, this threshold was originally reported for a complete computation, whereas we apply it here for two reasons: first, to enable direct comparison with~\cite{tourdegross}, and second, because inter-module measurements account for approximately $97\%$ of the total error in our model.

For a physical error rate of $p=10^{-3}$, the failure-probability threshold is reached at approximately $n=31$, $33$, and $34$ for $k=3$, $4$, and $5$, respectively.  For a physical error rate of $p=10^{-4}$, the projected circuit size required to reach the failure-probability threshold lies far beyond the range considered in our experiments $n \leq 450$, indicating that the threshold occurs only at substantially larger values of $n$.
% \begin{figure*}[t]
%     \centering  

%     \includegraphics[width=0.32\textwidth]%
%     {Figures/error_estimations/results_gross_1e-3_3_total_error.pdf}%
%     \hfill
%     \includegraphics[width=0.32\textwidth]%
%     {Figures/error_estimations/results_gross_1e-3_4_total_error.pdf}%
%     \hfill
%     \includegraphics[width=0.32\textwidth]%
%     {Figures/error_estimations/results_gross_1e-3_5_total_error.pdf}
%     \makebox[0.32\textwidth]{(a) $k=3$}\hfill
%     \makebox[0.32\textwidth]{(b) $k=3$}\hfill
%     \makebox[0.32\textwidth]{(c) $k=5$}

%     \includegraphics[width=0.32\textwidth]%
%     {Figures/error_estimations/results_gross_1e-3_3_failure_probability.pdf}%
%     \hfill
%     \includegraphics[width=0.32\textwidth]%
%     {Figures/error_estimations/results_gross_1e-3_4_failure_probability.pdf}%
%     \hfill
%     \includegraphics[width=0.32\textwidth]%
%     {Figures/error_estimations/results_gross_1e-3_5_failure_probability.pdf}
%     \makebox[0.32\textwidth]{(a) $k=3$}\hfill
%     \makebox[0.32\textwidth]{(b) $k=3$}\hfill
%     \makebox[0.32\textwidth]{(c) $k=5$}
%     \caption{Comparisons for inter-module movement for different factory counts}
%     \label{fig:saturation_curves}
% \end{figure*}





% \begin{table*}
%   \caption{Some Typical Commands}
%   \label{tab:commands}
%   \begin{tabular}{ccl}
%     \toprule
%     Command &A Number & Comments\\
%     \midrule
%     \texttt{{\char'134}author} & 100& Author \\
%     \texttt{{\char'134}table}& 300 & For tables\\
%     \texttt{{\char'134}table*}& 400& For wider tables\\
%     \bottomrule
%   \end{tabular}
% \end{table*}




% Conclusion
\section{Conclusion and Future Works}
\label{sec:con}
Multi-controlled Toffoli gates provide an important benchmark for quantum circuit compilation and resource optimization. In this work, we study their mapping onto BB-Code based fault-tolerant architectures and introduce a structure-aware placement strategy that exploits the binary-tree structure of optimal-depth MCT decompositions. Our results demonstrate that decomposition-aware placement can reduce inter-module communication and improve routing efficiency. More broadly, our work highlights the shift in compilation objectives from the NISQ regime, where circuits are mapped directly onto physical qubits, to the fault-tolerant regime, where logical modules, communication constraints, and magic-state factories are central for error-safe circuit execution. Extending this approach to other arithmetic and cryptographic applications, as well as to other MCT decompositions, fault-tolerant topologies, and resource models, provides a natural direction for future research.

%%% ACKNOWLEDGMENT
% \begin{acks}
% To be added later.
% \end{acks}


%\section*{Ethics and Privacy Statement}
% This section of your ACM work should discuss the potential societal risks that might result from its publication; two to three sentences related to the findings of your study, or new advancements made possible by their developed methods. The privacy and ethics statement should clearly address the broader impacts of their work as it relates to the authors' interpretation of privacy, fairness, safety, human rights, data sovereignty, or future misuse and any benefit/risk trade-off resulting from this research. We acknowledge that some papers may have minimal societal risks beyond those considered by institutional review boards, and the dimensions considered by any review of the user study design or dataset licenses could be provided in this statement.

%%% BIBLIOGRAPHY
% \newpage
\bibliographystyle{ACM-Reference-Format}
\bibliography{Reference}


%%% APPENDIX
\appendix

%\section{Appendix 1}

%\section{Appendix 2}



\end{document}
